\chapter{Mean Reversion and the Ornstein--Uhlenbeck Process}
\label{ch:mean_reversion_ou}
\chaptermark{Mean reversion and OU}

You have a time series that wanders around a constant value:
strays high, comes back; strays low, comes back. Can you trade
it? The naive answer is short the highs and buy the lows, but
turning ``it looks mean-reverting'' into a strategy requires
three quantitative answers in order:
\begin{enumerate}
  \item Is the series \emph{actually} mean-reverting, or does it
  only look that way? Most series that visually appear to revert
  are random walks that happened to visit the same level twice.
  Trading them on a mean-reversion thesis loses money.
  \item If the series \emph{is} mean-reverting, what is the
  half-life --- how long does the average excursion take to fade?
  A half-life of one minute is tradeable on a one-day horizon; a
  half-life of six months is not.
  \item Given the half-life, when do you enter, exit, and how do
  you size? The standard answer is the \textbf{$z$-score rule}:
  enter when the standardised deviation is large enough that the
  expected reversion is worth the capital, exit when it has
  decayed close to zero.
\end{enumerate}
This chapter answers all three with the
\textbf{Ornstein--Uhlenbeck (OU) process}, the canonical
continuous-time model of a mean-reverting variable. The OU model
solves (2) with an explicit half-life formula in one parameter;
it solves (3) by turning that parameter into entry and exit
thresholds; and it gives partial answers to (1) via three
statistical tests (Augmented Dickey--Fuller, variance ratio, Hurst
exponent). We follow Chan \citeChan{} and the more rigorous
Cartea--Jaimungal--Penalva \citeCartea{}, with one worked example
carried through every section.

Context. The Glosten--Milgrom posterior of \cref{ch:kyle_gm} is
a \emph{martingale}: no reversion, best prediction is the current
value. OU is the opposite. Efficient prices of tradeable
securities are martingales by no-arbitrage; \emph{spreads},
ETF \emph{premia}, and \emph{deviations} from a slow fair value
are mean-reverting. Recognising which is which is half the work;
the rest is the OU machinery.

\begin{backgroundbox}[Prerequisites]
From earlier in the book:
\begin{itemize}
  \item \Cref{ch:spreads_roll}: the Roll (1984) decomposition
  $\Cov(\Delta p_t, \Delta p_{t+1}) = -s^2/4$. Roll is the
  discrete-time trade-horizon mean-reverter; OU is the
  continuous-time generalisation that separates reversion speed
  from size.
  \item \Cref{ch:kyle_gm}: the Glosten--Milgrom posterior is a
  \emph{martingale}. Opposite worldview --- contrast deliberately.
  \item \Cref{ch:order_flow}: Hawkes processes modelled clustered
  \emph{arrivals}. OU is the mean-reversion analogue for
  clustered \emph{values} (drifting back to a mean rather than
  wandering off).
  \item \Cref{ch:market_making}: the Avellaneda--Stoikov
  $\volat$ for mid-price volatility is the same symbol and units
  as the OU driving-Brownian volatility here.
\end{itemize}
From outside the book:
\begin{itemize}
  \item Stochastic calculus, informally: Brownian motion $W_t$,
  It\^o integral $\int_0^t f(s)\, dW_s$, It\^o isometry
  $\E[(\int_0^t f\, dW_s)^2] = \int_0^t f^2\, ds$. Full
  machinery is built in \cref{ch:black_scholes}.
  \item OLS linear regression; the $t$-statistic, AR(1), and
  standard normal.
\end{itemize}
\end{backgroundbox}

\section{The Ornstein--Uhlenbeck process}
\label{sec:ch9_ou}

Introduced in 1930 by Ornstein and Uhlenbeck for a particle
subject to friction and thermal kicks, OU is the workhorse of
stochastic mean reversion in finance: every short-rate model in
\cref{ch:rates_derivatives} uses an OU-like SDE, every
cointegrated pair in \cref{ch:pairs_cointegration} treats its
residual as OU, and every mean-reversion strategy here is
parameterised by the OU half-life. (Symbols: $\theta$
mean-reversion rate, $\mu$ long-run mean, $\volat$ instantaneous
volatility, $W_t$ Brownian motion, $\tau_{1/2}$ half-life, $z$
standardised deviation.)

\begin{definition}[Ornstein--Uhlenbeck process]
\index{Ornstein--Uhlenbeck (OU) process|textbf}\index{OU process|textbf}\index{mean reversion!OU process|textbf}
\label{def:ou}
A stochastic process $(X_t)_{t \geq 0}$ taking values in $\R$ is
an \emph{Ornstein--Uhlenbeck process} with parameters
$(\theta, \mu, \volat)$ if it satisfies the stochastic
differential equation
\begin{equation}
  dX_t \;=\; \theta\,(\mu - X_t)\,dt \;+\; \volat\,dW_t,
  \qquad \theta > 0,\ \mu \in \R,\ \volat > 0,
  \label{eq:ou_sde}
\end{equation}
where $W_t$ is a standard Brownian motion. The parameter $\theta$
is the \emph{mean-reversion speed}, $\mu$ is the \emph{long-run
mean}, and $\volat$ is the \emph{instantaneous volatility}.
\end{definition}

The drift $\theta(\mu - X_t)$ does the work. When $X_t > \mu$ the
drift is negative and pulls down; when $X_t < \mu$ it pulls up,
with strength proportional to the distance from $\mu$. Larger
$\theta$ means faster reversion. The noise $\volat\, dW_t$ kicks
$X_t$ around as it reverts; without it $X_t$ would decay
exponentially toward $\mu$.

\begin{figure}[ht]
\centering
\begin{tikzpicture}
\begin{axis}[
  width=10cm, height=5cm,
  xlabel={Time $t$},
  ylabel={$X_t$},
  xmin=0, xmax=10,
  ymin=-2.8, ymax=2.8,
  ytick={-2,0,2},
  xtick={0,2,4,6,8,10},
  grid=major, grid style={dashed,gray!25},
  title={Ornstein--Uhlenbeck sample path ($\mu=0$)},
  title style={font=\small\bfseries},
]
\addplot[dashed, black!45, thick]
  coordinates {(0,0)(10,0)} node[right, font=\small]{$\mu$};
\addplot[blue!70, thick, smooth] coordinates {
  (0,1.5)(0.5,1.8)(1.0,1.3)(1.5,0.6)(2.0,0.1)(2.5,-0.4)
  (3.0,-0.9)(3.5,-1.5)(4.0,-1.9)(4.5,-1.4)(5.0,-0.7)
  (5.5,-0.1)(6.0,0.5)(6.5,1.1)(7.0,1.7)(7.5,2.1)
  (8.0,1.5)(8.5,0.9)(9.0,0.3)(9.5,-0.3)(10.0,0.1)
};
\draw[-Stealth, red!65, thick]
  (axis cs:0.5,1.8) -- (axis cs:0.5,0.4)
  node[right, font=\tiny, red!70]{pull};
\draw[-Stealth, red!65, thick]
  (axis cs:4.0,-1.9) -- (axis cs:4.0,-0.5)
  node[right, font=\tiny, red!70]{pull};
\draw[-Stealth, red!65, thick]
  (axis cs:7.5,2.1) -- (axis cs:7.5,0.6)
  node[right, font=\tiny, red!70]{pull};
\end{axis}
\end{tikzpicture}
\caption{Simulated Ornstein--Uhlenbeck path. Red arrows show the
mean-reversion force: whenever $X_t$ strays from $\mu = 0$,
the drift term $\theta(\mu - X_t)$ pulls it back.
The half-life $\tau_{1/2} = \log(2)/\theta$ sets the typical
return-to-mean timescale. Larger $\theta$ produces faster, tighter
oscillations; $\theta \to 0$ recovers a random walk.}
\label{fig:ou_path}
\end{figure}

\subsection{Solving the SDE}
\label{subsec:ch9_solve}

Equation~\eqref{eq:ou_sde} is a first-order linear SDE and solves
in closed form by the integrating-factor trick. Multiply both
sides by $e^{\theta t}$:
\begin{equation}
  e^{\theta t}\,dX_t \;+\; \theta\,e^{\theta t}\,X_t\,dt
  \;=\; \theta\,e^{\theta t}\,\mu\,dt \;+\; \volat\,e^{\theta t}\,dW_t.
  \label{eq:ou_step1}
\end{equation}
The left-hand side is a perfect differential. Applying It\^o's
product rule to $f(t, X_t) = e^{\theta t} X_t$ (linear in $X_t$,
so the second-order term vanishes),
\begin{align}
  d\!\left(e^{\theta t} X_t\right)
  \;&=\; \theta\,e^{\theta t}\,X_t\,dt
       \;+\; e^{\theta t}\,dX_t \nonumber \\[0.25em]
  \;&=\; e^{\theta t}\,dX_t \;+\; \theta\,e^{\theta t}\,X_t\,dt,
  \label{eq:ou_step2}
\end{align}
which is exactly the left-hand side of \eqref{eq:ou_step1}. So
\eqref{eq:ou_step1} reads
\begin{equation}
  d\!\left(e^{\theta t} X_t\right)
  \;=\; \theta\,\mu\,e^{\theta t}\,dt
  \;+\; \volat\,e^{\theta t}\,dW_t.
  \label{eq:ou_step3}
\end{equation}
Integrate both sides from $0$ to $t$:
\begin{align}
  e^{\theta t}\,X_t \;-\; X_0
  \;&=\; \theta\,\mu\,\int_0^t e^{\theta s}\,ds
       \;+\; \volat\,\int_0^t e^{\theta s}\,dW_s \nonumber \\[0.25em]
  \;&=\; \mu\,\bigl(e^{\theta t} - 1\bigr)
       \;+\; \volat\,\int_0^t e^{\theta s}\,dW_s.
  \label{eq:ou_step4}
\end{align}
Divide through by $e^{\theta t}$ and rearrange:
\begin{equation}
  \boxed{\;
  X_t \;=\; X_0\,e^{-\theta t}
       \;+\; \mu\,\bigl(1 - e^{-\theta t}\bigr)
       \;+\; \volat\,\int_0^t e^{-\theta(t-s)}\,dW_s.
  \;}
  \label{eq:ou_solution}
\end{equation}
Three pieces. $X_0 e^{-\theta t}$ is the \emph{exponential decay
of the initial condition}: the influence of $X_0$ shrinks at rate
$\theta$. $\mu(1 - e^{-\theta t})$ is the \emph{exponential growth
of the mean toward $\mu$}: the deterministic part heads to $\mu$
regardless of start. The It\^o integral is the \emph{accumulated
noise}, exponentially weighted so recent shocks count more --- old
kicks have already reverted away.

\subsection{Mean and variance at finite time}
\label{subsec:ch9_meanvar}

Take expectations of \eqref{eq:ou_solution}. The It\^o integral
of a deterministic integrand is a zero-mean Gaussian, so it drops
out of the expectation:
\begin{equation}
  \E[X_t]
  \;=\; X_0\,e^{-\theta t} \;+\; \mu\,\bigl(1 - e^{-\theta t}\bigr).
  \label{eq:ou_mean}
\end{equation}
The expected value moves exponentially from $X_0$ toward $\mu$ at
rate $\theta$.

For the variance, the deterministic terms contribute zero. The
stochastic piece $\volat \int_0^t e^{-\theta(t-s)}\, dW_s$ yields
by It\^o isometry
\begin{align}
  \Var\!\left[\volat\,\int_0^t e^{-\theta(t-s)}\,dW_s\right]
  \;&=\; \volat^2\,\int_0^t e^{-2\theta(t-s)}\,ds.
  \label{eq:ou_var_step1}
\end{align}
Substitute $u = t - s$ and flip the limits:
\begin{align}
  \int_0^t e^{-2\theta(t-s)}\,ds
  \;&=\; \int_0^t e^{-2\theta u}\,du
  \;=\; \frac{1 - e^{-2\theta t}}{2\theta}. \label{eq:ou_var_step2}
\end{align}
Therefore
\begin{equation}
  \Var[X_t] \;=\; \frac{\volat^2}{2\theta}\,
                 \bigl(1 - e^{-2\theta t}\bigr).
  \label{eq:ou_var_finite}
\end{equation}
At $t = 0$ this is zero; as $t \to \infty$ it grows monotonically
to a finite limit:
\begin{equation}
  \Var[X_t] \;\xrightarrow{t \to \infty}\;
  \sigma_\infty^2 \;\defeq\; \frac{\volat^2}{2\theta}.
  \label{eq:ou_var_stationary}
\end{equation}
The finite limit is why a mean-reverter can be \emph{traded}.
A random walk has variance growing linearly forever, so ``it
will come back'' has unbounded downside. OU has \emph{bounded}
variance --- values stay in a tight neighbourhood of $\mu$ with
a finite typical excursion size.

\subsection{The stationary distribution}
\label{subsec:ch9_stationary}

For each finite $t$, $X_t$ in \eqref{eq:ou_solution} is Gaussian
(linear combination of Gaussian initial condition and Gaussian
It\^o integral), with mean and variance
\eqref{eq:ou_mean}--\eqref{eq:ou_var_finite}. As $t \to \infty$
both converge: mean to $\mu$, variance to
$\sigma_\infty^2 = \volat^2/(2\theta)$. The limiting distribution
is
\begin{equation}
  X_\infty \;\sim\; \Normal\!\left(\mu,\;\frac{\volat^2}{2\theta}\right).
  \label{eq:ou_stationary_dist}
\end{equation}
This is the \emph{stationary distribution}: if $X_0$ is drawn
from it, $X_t$ has the same distribution for every $t \geq 0$.
The process never settles to a single point --- noise keeps
kicking it --- but it settles to a stable Gaussian envelope of
typical width $\sigma_\infty = \volat/\sqrt{2\theta}$ around
$\mu$.

\subsection{Autocorrelation}
\label{subsec:ch9_autocorr}

Two snapshots of the OU process at times $t$ and $t + h$ are
correlated; the correlation decays exponentially with the lag
$h$. In the stationary regime,
\begin{equation}
  \Corr(X_t,\, X_{t+h}) \;=\; e^{-\theta h}.
  \label{eq:ou_autocorr}
\end{equation}
The derivation starts from $\Cov(X_t, X_{t+h}) =
\Cov(X_t, X_t e^{-\theta h} + \text{noise}) = e^{-\theta h}
\sigma_\infty^2$ and divides by $\sigma_\infty^2$; Cartea
\citeCartea{} (Ch.~5) reproduces the steps. The key fact: the
autocorrelation of a stationary OU process is a \emph{single
decaying exponential} with the same rate $\theta$ as the drift.
The OU process has one timescale, $1/\theta$.

\subsection{Half-life}
\label{subsec:ch9_halflife}

The mean of $X_t$, given a starting point $X_0$, decays toward
$\mu$ as $\E[X_t] - \mu = (X_0 - \mu)\,e^{-\theta t}$ from
\eqref{eq:ou_mean}. The \emph{half-life} of the process is the
time it takes for the expected deviation from $\mu$ to shrink to
half its initial size:
\begin{equation}
  e^{-\theta\,h_{1/2}} \;=\; \tfrac{1}{2}
  \quad\Longleftrightarrow\quad
  \theta\,h_{1/2} \;=\; \ln 2
  \quad\Longleftrightarrow\quad
  h_{1/2} \;=\; \frac{\ln 2}{\theta}.
  \label{eq:ou_halflife}
\end{equation}
The half-life is the single most important quantity for a
trader: in the natural units of the series, how long on average
you must wait for an excursion to revert. Everything else ---
entry threshold, exit threshold, trade duration, capital
allocation --- is downstream.

\begin{keyfactbox}[OU process summary]
The Ornstein--Uhlenbeck SDE
\(dX_t = \theta(\mu - X_t)\,dt + \volat\,dW_t\) with
$\theta > 0$ has the following properties.
\begin{itemize}
  \item \textbf{Closed-form solution}
  $X_t = X_0 e^{-\theta t} + \mu(1 - e^{-\theta t})
       + \volat\int_0^t e^{-\theta(t-s)}\,dW_s$.
  \item \textbf{Stationary distribution}
  $X_\infty \sim \Normal(\mu,\, \volat^2/(2\theta))$.
  Stationary standard deviation
  $\sigma_\infty = \volat/\sqrt{2\theta}$.
  \item \textbf{Autocorrelation}
  $\Corr(X_t, X_{t+h}) = e^{-\theta h}$. Single exponential
  timescale $1/\theta$.
  \item \textbf{Half-life}
  $h_{1/2} = (\ln 2)/\theta$. Time for an expected
  excursion to fade by half.
\end{itemize}
\end{keyfactbox}

\subsection{A worked numerical example}
\label{subsec:ch9_worked_ou}

Take $\mu = 100$, $\theta = 0.3$, $\volat = 2$, $X_0 = 105$
(five units above the mean). Time units are arbitrary --- think
``hours'' or ``Prosperity ticks''.

\paragraph{Half-life.} From \eqref{eq:ou_halflife},
\begin{equation}
  h_{1/2} \;=\; \frac{\ln 2}{0.3}
            \;\approx\; \frac{0.693147}{0.3}
            \;\approx\; 2.310 \text{ time units}.
  \label{eq:ch9_worked_halflife}
\end{equation}
A deviation of size $5$ above the mean is expected to shrink to
$2.5$ in about $2.31$ time units, to $1.25$ after another $2.31$,
and so on.

\paragraph{Stationary spread.} The stationary standard deviation
is
\begin{equation}
  \sigma_\infty \;=\; \frac{\volat}{\sqrt{2\theta}}
                \;=\; \frac{2}{\sqrt{0.6}}
                \;=\; \frac{2}{0.7746\ldots}
                \;\approx\; 2.582.
  \label{eq:ch9_worked_sigmainf}
\end{equation}
So $X_t$ stays mostly within $\mu \pm 2\sigma_\infty
\approx [94.84, 105.16]$ ($95\%$ mass by the Gaussian rule).

\paragraph{Expected trajectory and variance build-up.} From
\eqref{eq:ou_mean} and \eqref{eq:ou_var_finite}, plug in:
\begin{center}
\begin{tabular}{rrr}
\toprule
$t$ & $\E[X_t]$ & $\Var[X_t]$ \\
\midrule
$0$        & $105.000$ & $0.000$ \\
$1$        & $103.704$ & $3.008$ \\
$2$        & $102.744$ & $4.659$ \\
$h_{1/2}=2.310$ & $102.500$ & $5.000$ \\
$5$        & $101.116$ & $6.335$ \\
$10$       & $100.249$ & $6.650$ \\
$20$       & $100.012$ & $6.667$ \\
$\infty$   & $100.000$ & $6.667$ \\
\bottomrule
\end{tabular}
\end{center}
Sanity checks: at $t = h_{1/2}$ the expected deviation is
$2.5 = (105-100)/2$, as the half-life formula demands; at
$t = \infty$ the variance is $\volat^2/(2\theta) = 4/0.6 = 6.667$,
matching \eqref{eq:ou_var_stationary}, with
$\sqrt{6.667} \approx 2.582$, matching
\eqref{eq:ch9_worked_sigmainf}.

\begin{intuitionbox}
The half-life is the only number that matters. If it is one day
and your holding window is five days, you live through several
revert-and-reset cycles --- profitable in expectation. If it is
six months and your window is five days, no excursion reverts in
time; the strategy becomes a directional bet. Half-life $\ll$
horizon is the gate that separates ``trade this'' from ``leave
it alone''. Everything else is calibration.
\end{intuitionbox}

\section{Testing for mean reversion}
\label{sec:ch9_tests}

Most financial time series are \emph{not} OU --- they are random
walks. Equities, FX rates, and commodity futures are
martingale-like at any human-tradeable horizon. Trading a random
walk on a mean-reversion thesis loses money: there is no
reversion to trade against.

Before applying OU machinery you must \emph{reject the
random-walk null}. Three tests are standard and a sell-side
stat-arb desk runs all three: they fail in different ways, so a
series surviving all three is much more credible than one
surviving only one.

\subsection{The Augmented Dickey--Fuller test}
\label{subsec:ch9_adf}

The \textbf{Augmented Dickey--Fuller (ADF) test} is the canonical
test of stationarity in a univariate time series. In plain terms
it asks: ``is there a statistically significant force pulling
this series back toward a constant level, or does it wander
freely?'' The null is no such force (a random walk, technically a
\emph{unit root}); rejecting the null is evidence of mean
reversion. It runs the regression
\begin{equation}
  \Delta X_t \;=\; \alpha
                  \;+\; \beta\,X_{t-1}
                  \;+\; \sum_{i=1}^p \gamma_i\,\Delta X_{t-i}
                  \;+\; \epsilon_t,
  \label{eq:adf_regression}
\end{equation}
where $\Delta X_t \defeq X_t - X_{t-1}$ is the first difference,
$\alpha$ an intercept (drift), and the $p$ lagged differences are
nuisance terms that soak up serial correlation in $\epsilon_t$.
The null is $H_0\colon \beta = 0$: under $H_0$ the level
$X_{t-1}$ has no predictive power for $\Delta X_t$ --- $X_t$ is a
random walk. Under $H_1\colon \beta < 0$, a high $X_{t-1}$ makes
$\Delta X_t$ more negative on average --- mean reversion in
difference form.

The ADF test statistic is the $t$-statistic of $\hat\beta$. Under
the null its distribution is \emph{not} Student's $t$: the
regressor $X_{t-1}$ is itself a random walk and standard
regression theory breaks. Correct critical values are tabulated
by MacKinnon (cited via \citeChan); a typical $5\%$ critical
value at long sample sizes is $\approx -2.86$, and the test
rejects when the observed $t$-statistic is more \emph{negative}.
Choose lag count $p$ by AIC or BIC; Chan recommends starting at
$p = 1$ and increasing until residual autocorrelation
disappears.

For the running example ($\theta = 0.3$, $\mu = 100$,
$\volat = 2$, $N = 1000$), ADF typically yields a statistic
around $-7$, comfortably past $-2.86$, and rejects the
random-walk null.

\begin{pitfallbox}[The ADF test has low power in short samples]
The ADF test is conservative: it fails to reject the random-walk
null much more often than it should when the true process is OU
but the sample is short or $\theta$ is small. With $N=200$ and
$\theta = 0.05$ (half-life $\approx 14$), simulated ADF rejects
at $5\%$ only $\sim 30\%$ of the time even though every
simulation came from an OU. A non-rejecting ADF is consistent
with both a random walk and a slow reverter. Corroborate with a
half-life estimate and a variance-ratio test; never trust a
single ADF reading on $< 250$ observations.
\end{pitfallbox}

\subsection{The variance ratio test}
\label{subsec:ch9_vrt}

The \textbf{variance ratio (VR) test} of Lo and MacKinlay uses a
simpler observation: under a random walk the variance of
$k$-period changes equals $k$ times the variance of one-period
changes, because increments are independent. Define
\begin{equation}
  VR(k) \;\defeq\; \frac{\Var(X_t - X_{t-k})}{k\,\Var(X_t - X_{t-1})}.
  \label{eq:vrt_def}
\end{equation}
Under the random-walk null $VR(k) = 1$ for every $k$. Under mean
reversion, negatively-correlated changes (a step up tends to be
followed by a step down) give $VR(k) < 1$; under trending,
$VR(k) > 1$.

For the running example, a direct calculation using the OU
autocorrelation \eqref{eq:ou_autocorr} gives, in the stationary
regime,
\begin{equation}
  VR(k) \;=\; \frac{1 - e^{-\theta k}}{k\,\bigl(1 - e^{-\theta}\bigr)},
  \label{eq:vrt_ou}
\end{equation}
which at $\theta = 0.3$ takes values
\begin{center}
\begin{tabular}{rl}
\toprule
$k$ & $VR(k)$ \\
\midrule
$1$  & $1.000$ \\
$2$  & $0.870$ \\
$5$  & $0.600$ \\
$10$ & $0.367$ \\
$20$ & $0.192$ \\
\bottomrule
\end{tabular}
\end{center}
The ratio falls toward zero as $k$ grows: long-horizon changes
are bounded by the OU range, while one-period changes are
dominated by noise. The Lo--MacKinlay test maps these ratios to
a $z$-statistic with known asymptotic distribution; for a
long-enough sample, rejection at any $k \in [5, 20]$ is almost
mechanical for a true OU with non-trivial $\theta$.

\subsection{The Hurst exponent}
\label{subsec:ch9_hurst}

The \textbf{Hurst exponent} $H$ is a third diagnostic, defined so
that the long-run scaling of the range of a process is
$R_t \propto t^H$. Three regimes:
\begin{itemize}
  \item $H = 0.5$ --- pure random walk.
  \item $H < 0.5$ --- mean-reverting (the range grows slower than
  $\sqrt{t}$, indicating that excursions get pulled back).
  \item $H > 0.5$ --- trending (the range grows faster than
  $\sqrt{t}$, indicating persistent moves).
\end{itemize}
The classical estimator is the rescaled-range ($R/S$) procedure
of Mandelbrot, reviewed in \citeChan{}; modern practice uses a
multifractal extension. $H$ returns one number that classifies
the regime, but it is the noisiest of the three tests --- Chan
recommends using it only as a tiebreaker after ADF and VR.

The three tests are not redundant: ADF asks ``can a linear
mean-reversion regression beat a random walk?'', VR asks ``do
multi-period variances accumulate sublinearly?'', Hurst asks
``does the range scale slower than $\sqrt{t}$?''. Different
summaries of the same series; a stat-arb shop runs all three
before committing capital. The Frankfurt Hedgehogs Prosperity 3
writeup (\texttt{imc-prosperity-3/README.md}) explicitly checked
for OU-style mean reversion in PICNIC\_BASKET premia and
VOLCANIC\_ROCK before turning on those strategies.

\section{Estimating OU parameters from data}
\label{sec:ch9_fit}

Once a series is classed as mean-reverting, estimate $\theta$,
$\mu$, $\volat$ from data. The standard trick: the OU SDE has an
exact discrete-time representation as an AR(1), whose parameters
fall out of a single OLS regression.

\subsection{The discrete-time AR(1) representation}
\label{subsec:ch9_ar1}

Sample $X_t$ at uniform intervals of length $\Delta t$ and apply
the closed-form solution \eqref{eq:ou_solution} between $t$ and
$t + \Delta t$:
\begin{equation}
  X_{t+\Delta t}
  \;=\; X_t\,e^{-\theta \Delta t}
       \;+\; \mu\,\bigl(1 - e^{-\theta \Delta t}\bigr)
       \;+\; \volat\,\int_t^{t+\Delta t} e^{-\theta(t+\Delta t - s)}\,dW_s.
  \label{eq:ou_discrete_step}
\end{equation}
The integral is a zero-mean Gaussian with variance (by the same
It\^o-isometry argument)
\begin{equation}
  \Var\!\left[\volat\,\int_t^{t+\Delta t} e^{-\theta(t+\Delta t - s)}\,dW_s\right]
  \;=\; \frac{\volat^2}{2\theta}\,\bigl(1 - e^{-2\theta\Delta t}\bigr).
  \label{eq:ou_eps_var}
\end{equation}
Define
\begin{equation}
  \phi \;\defeq\; e^{-\theta \Delta t}, \qquad
  a   \;\defeq\; \mu\,(1 - \phi), \qquad
  \sigma_\epsilon^2 \;\defeq\; \frac{\volat^2}{2\theta}\,\bigl(1 - \phi^2\bigr),
  \label{eq:ou_ar1_reparam}
\end{equation}
and the discrete-time recursion \eqref{eq:ou_discrete_step}
becomes
\begin{equation}
  X_{t+\Delta t} \;=\; \phi\,X_t \;+\; a \;+\; \epsilon_t,
  \qquad \epsilon_t \overset{\iid}{\sim}\Normal(0,\,\sigma_\epsilon^2).
  \label{eq:ou_ar1}
\end{equation}
An exact AR(1). Under uniform sampling the AR(1) is not an
approximation: the OU SDE and the AR(1) are literally the same
model.

\subsection{OLS recovery of \texorpdfstring{$\theta$, $\mu$, $\volat$}{theta, mu, sigma}}
\label{subsec:ch9_ols}

Regress $X_{t+\Delta t}$ on $X_t$ by OLS. This returns slope
$\hat\phi$, intercept $\hat a$, and residual std
$\hat\sigma_\epsilon$. Invert \eqref{eq:ou_ar1_reparam}:
\begin{align}
  \hat\theta \;&=\; -\frac{\ln \hat\phi}{\Delta t},
                                            \label{eq:ou_theta_hat} \\[0.25em]
  \hat\mu    \;&=\; \frac{\hat a}{1 - \hat\phi},
                                            \label{eq:ou_mu_hat} \\[0.25em]
  \hat\volat \;&=\; \hat\sigma_\epsilon\,
                  \sqrt{\frac{2\,\hat\theta}{1 - \hat\phi^2}}.
                                            \label{eq:ou_sigma_hat}
\end{align}
The first inverts $\phi = e^{-\theta \Delta t}$; the second
inverts $a = \mu(1-\phi)$; the third solves \eqref{eq:ou_eps_var}
for $\volat$.

\paragraph{Worked verification.} Simulate an OU trajectory with
$\theta = 0.3$, $\mu = 100$, $\volat = 2$, $\Delta t = 1$,
$N = 5000$; true $\phi = e^{-0.3} \approx 0.7408$. OLS returns
$\hat\phi \approx 0.734$, $\hat\theta \approx 0.310$,
$\hat\mu \approx 100.04$, $\hat\volat \approx 2.00$, within
sampling error. On $N = 500$ points the same procedure returns
$\hat\theta \in [0.2, 0.4]$ depending on seed --- $\theta$ is
the hardest of the three to pin down, and narrow confidence
intervals on the half-life require a lot of data.

\begin{keyfactbox}[OU fitting via AR(1)]
To fit $(\theta, \mu, \volat)$ to a uniformly sampled time
series $\{X_t\}$ at spacing $\Delta t$:
\begin{enumerate}
  \item Regress $X_{t+\Delta t}$ on $X_t$ by OLS, obtaining
  slope $\hat\phi$, intercept $\hat a$, residual standard
  deviation $\hat\sigma_\epsilon$.
  \item Recover $\hat\theta = -\ln(\hat\phi)/\Delta t$.
  \item Recover $\hat\mu = \hat a / (1 - \hat\phi)$.
  \item Recover
  $\hat\volat = \hat\sigma_\epsilon\,\sqrt{2\hat\theta/(1-\hat\phi^2)}$.
\end{enumerate}
The half-life follows immediately:
$\hat h_{1/2} = (\ln 2)/\hat\theta$.
\end{keyfactbox}

\section{The \texorpdfstring{$z$}{z}-score entry/exit rule}
\label{sec:ch9_zscore}

With $(\hat\theta, \hat\mu, \hat\volat)$ in hand, define the
\emph{$z$-score} of $X_t$:
\begin{equation}
  z_t \;\defeq\; \frac{X_t - \hat\mu}{\hat\sigma_\infty},
  \qquad
  \hat\sigma_\infty \;=\; \frac{\hat\volat}{\sqrt{2\hat\theta}}.
  \label{eq:zscore_def}
\end{equation}
The $z$-score is the current deviation from the long-run mean
in stationary-std units. Under OU in steady state
$z_t \sim \Normal(0, 1)$, so ``$z = 2$'' is a $\sim 2.5\%$ tail
event. The threshold $z_{\text{entry}} = 2$ is therefore picked
not arbitrarily but because a stationary OU process spends only
$\sim 5\%$ of its time outside $\pm 2\sigma_\infty$: when it does
stray that far, reversion is highly probable and the expected
move back toward $\mu$ is large enough to pay for transaction
costs.

The standard \textbf{$z$-score trading rule}, used by
practitioners since the 1980s and codified by Chan
\citeChan{}, is:
\begin{itemize}
  \item \textbf{Enter short}: when $z_t > z_{\text{entry}}$
  (e.g.\ $z_{\text{entry}} = 2$). The series is high relative
  to its long-run mean and is expected to fall back.
  \item \textbf{Enter long}: when $z_t < -z_{\text{entry}}$.
  Symmetrically.
  \item \textbf{Exit}: when $|z_t| < z_{\text{exit}}$ (e.g.\
  $z_{\text{exit}} = 0.5$). The deviation has decayed back to
  near zero and there is no further reversion to harvest.
  \item \textbf{Stop out}: when $|z_t| > z_{\text{stop}}$ (e.g.\
  $z_{\text{stop}} = 4$). Four standard deviations from the
  supposed mean means the OU model is wrong, the parameters are
  stale, or the regime has changed. Exit at a controlled loss
  before the position blows up.
\end{itemize}
The asymmetry between $z_{\text{entry}}$ and $z_{\text{exit}}$
is deliberate: requiring $|z_t| < z_{\text{entry}}$ to exit would
never exit on a deviation that decayed to the entry level and
bounced. The slack between $z_{\text{exit}}$ and
$z_{\text{entry}}$ ensures most trades close on a meaningful
reversion.

\subsection{Expected trade duration}
\label{subsec:ch9_duration}

How long does the typical trade last? Use the OU mean trajectory
$\E[X_t] - \mu = (X_0 - \mu)e^{-\theta t}$. In $z$-units, an
excursion of size $z_{\text{entry}}$ decays in expectation to
$z_{\text{exit}}$ after time $h$:
\begin{equation}
  z_{\text{entry}}\,e^{-\theta h} \;=\; z_{\text{exit}}
  \;\;\Longrightarrow\;\;
  h \;=\; \frac{1}{\theta}\,\ln\!\left(\frac{z_{\text{entry}}}{z_{\text{exit}}}\right).
  \label{eq:zscore_duration}
\end{equation}
The expected trade duration. With $z_{\text{entry}} = 2$,
$z_{\text{exit}} = 0.5$, $\theta = 0.3$,
\begin{equation}
  h \;=\; \frac{\ln 4}{0.3}
       \;\approx\; \frac{1.3863}{0.3}
       \;\approx\; 4.621 \text{ time units}.
  \label{eq:ch9_worked_duration}
\end{equation}
Sanity check: about twice the half-life
$h_{1/2} \approx 2.310$, as expected --- two halvings
($z = 2 \to 1 \to 0.5$) rather than one.

Equation~\eqref{eq:zscore_duration} is a
deterministic-trajectory estimate, ignoring the noise around the
mean trajectory. The exact mean first-passage time involves an
integral over the OU generator with no closed form for general
thresholds; for threshold-setting and Kelly sizing
(\cref{ch:kelly_risk}) the deterministic approximation is
adequate.

\begin{keyfactbox}[The $z$-score OU strategy]
Given OU parameter estimates $(\hat\theta, \hat\mu, \hat\volat)$
and stationary std $\hat\sigma_\infty = \hat\volat/\sqrt{2\hat\theta}$:
\begin{itemize}
  \item Compute the live $z$-score
  $z_t = (X_t - \hat\mu)/\hat\sigma_\infty$.
  \item \textbf{Enter} a position of size
  $-\operatorname{sgn}(z_t)$ when $|z_t| > z_{\text{entry}}$.
  \item \textbf{Exit} when $|z_t| < z_{\text{exit}}$.
  \item \textbf{Stop out} when $|z_t| > z_{\text{stop}}$.
  \item \textbf{Expected duration} of a trade is
  $h \approx \theta^{-1}\,\ln(z_{\text{entry}}/z_{\text{exit}})$.
\end{itemize}
Refit $(\hat\theta, \hat\mu, \hat\volat)$ on a rolling window
(see \cref{sec:ch9_walkforward}) so the parameters track the
current regime.
\end{keyfactbox}

\begin{intuitionbox}
A $z$-score entry is a bet that the half-life is shorter than
your deadline. With trading horizon $H$ and half-life $h_{1/2}$,
you get $H/h_{1/2}$ revert-and-reset cycles in expectation. Ratio
$\gg 1$: deviations come back inside the window and the strategy
harvests them. Ratio $\ll 1$: you mark to market for many
half-lives without reversion --- functionally a directional bet.
The thresholds set how far from the mean you wait before
entering, but they do not change the half-life. Pick the asset
first ($h_{1/2} < H$), then the thresholds.
\end{intuitionbox}

\section{Walk-forward estimation}
\label{sec:ch9_walkforward}

OU parameters drift. The mean of a stat-arb residual shifts with
the company's business; the mean-reversion speed shifts with
regime; the volatility shifts continuously. A model fit in
January is substantially wrong by July, and a model fit on the
entire historical record uses post-trade information to compute
pre-trade $z$-scores --- look-ahead bias that inflates simulated
profits and disappears live.

The fix is \textbf{walk-forward estimation}. Pick a fitting
window $W$ and refit cadence $R$. At each time $t$:
\begin{enumerate}
  \item Fit the OU model on the most recent $W$ observations
  $X_{t-W}, \ldots, X_{t-1}$, using only data \emph{strictly
  before} time $t$.
  \item Compute the live $z$-score
  $z_t = (X_t - \hat\mu_t)/\hat\sigma_{\infty,t}$ from the
  parameters fit on the prior window.
  \item Generate the trading signal from $z_t$.
  \item Every $R$ steps, repeat the fit on the latest window.
\end{enumerate}
Typical Prosperity-scale numbers: $W \in [500, 5000]$,
$R \in [50, 1000]$, scaled by half-life (shorter half-life
$\Rightarrow$ shorter $W$ and $R$). On a sell-side single-name
desk, $W$ is one to two trading years and $R$ is one trading
day.

\begin{pitfallbox}[Look-ahead bias in OU fitting]
The most common mistake on a first OU strategy: fit $\hat\mu$ on
the entire dataset and use it for every historical $z$-score.
This feeds \emph{future} data into the strategy. Simulated
$z$-scores are systematically better-centred than in live
trading, entries are timed more cleverly than possible, and the
backtest shows fictional profit. Same trap for $\hat\theta$ and
$\hat\volat$. The cure: parameter values used at time $t$ must
depend only on $X_s$ for $s < t$. \Cref{ch:backtesting} returns
to look-ahead bias as a category.
\end{pitfallbox}

\section{A short reference implementation}
\label{sec:ch9_code}

The OU fitting pipeline fits in under two dozen lines of Python.
The reference below is a \emph{pure function} from a numpy array
of past observations (with optional sampling interval) to
$(\hat\theta, \hat\mu, \hat\volat)$. The signal generator on top
is one line per the keyfactbox in \cref{sec:ch9_zscore}.

\begin{lstlisting}[caption={OU fitting via AR(1) as a pure function.}]
import numpy as np

def fit_ou(x: np.ndarray,
           dt: float = 1.0) -> tuple[float, float, float]:
    """Fit OU parameters (theta, mu, sigma) via the AR(1) mapping."""
    y, xp = x[1:], x[:-1]
    # np.polyfit returns highest-order coef first => (slope, intercept)
    phi, a = np.polyfit(xp, y, 1)
    theta = -np.log(phi) / dt
    mu    = a / (1.0 - phi)
    resid = y - (phi * xp + a)
    var_eps = float(np.var(resid, ddof=2))
    sigma = float(np.sqrt(2.0 * theta * var_eps / (1.0 - phi**2)))
    return float(theta), float(mu), sigma
\end{lstlisting}

Three production remarks not encoded above. First, the function
assumes $\hat\phi \in (0, 1)$; $\hat\phi \leq 0$ does not satisfy
OU and $-\ln(\hat\phi)$ is undefined, so the caller should guard
and refuse to trade. Second, \texttt{ddof=2} accounts for the
two OLS-estimated parameters $(\phi, a)$; without it
$\hat\volat$ is slightly biased downward on small samples.
Third, this function is meant to be called inside a walk-forward
loop --- calling it once on full history is the look-ahead trap
above.

\begin{prosperitybox}[title={Prosperity example: KELP}]
KELP is the canonical mildly-mean-reverting Prosperity product.
Per Frankfurt Hedgehogs
(\texttt{imc-prosperity-3/README.md}), KELP's true price follows
a slow random walk with a small mean-reverting overlay from
integer-tick rounding of an internal floating-point fair value:
``a minor technical edge stemming from the fact that the true
price was internally a floating-point value and orders could
only be posted at integer levels (creating slight mean-reversion
tendencies after ticks).'' KELP is the ideal teaching example
for OU fitting on Prosperity data.

\textbf{Procedure.}
\begin{enumerate}
  \item Pull the KELP wall-mid
  (\cref{ch:price_value}) from the first $1000$ ticks of a round
  --- the fitting window $W$.
  \item Run \texttt{fit\_ou} from \cref{sec:ch9_code} with
  $\Delta t = 1$.
  \item Compute $h_{1/2} = (\ln 2)/\hat\theta$ and
  $\hat\sigma_\infty = \hat\volat/\sqrt{2\hat\theta}$.
  \item Typically $\hat\theta \in [0.02, 0.05]$, giving a
  half-life of $15$--$30$ ticks.
  \item Set $z_{\text{entry}} = 1.5$, $z_{\text{exit}} = 0.3$
  (tighter than the section default because fills are quick and
  waiting for $|z| > 2$ misses trades). Expected duration
  $h = \ln(5)/\hat\theta \approx 1.61/\hat\theta$, so $30$--$80$
  ticks per trade. Over a $10\,000$-tick episode, comfortably
  more than $50$ trades.
  \item Refit every $250$ ticks on a rolling $1000$-tick window.
  $\hat\mu$ tracks the slow random walk; $\hat\theta$ stays
  approximately constant (the rounding mechanism does not change).
\end{enumerate}

\textbf{What the top teams actually did.} The Frankfurt
Hedgehogs' final KELP handler
(\texttt{imc-prosperity-3/FrankfurtHedgehogs\_polished.py},
\texttt{DynamicTrader.get\_orders()}) does \emph{not} run a
literal $z$-score OU strategy on KELP. It quotes one tick inside
the wall bid/ask --- the fixed-offset rule used on
RAINFOREST\_RESIN in \cref{ch:market_making} --- and lets the
wall-mid update absorb the slow random walk. They treated KELP
as ``unpredictable apart from a tiny rounding effect'' and
concluded the effect was too small to extract reliably
($\sim 5\,000$ seashells per round, mostly market making). The
lesson: OU is the right tool to \emph{diagnose} that KELP's
reversion is too weak to trade as a pure stat-arb signal, and
the right tool to \emph{trade} a stronger mean-reverter like the
PICNIC\_BASKET premium of \cref{ch:basket_etf_arb}, where the
half-life is minutes and $|z| > 2$ happens regularly.
\end{prosperitybox}

\begin{gsbox}
On a sell-side equity stat-arb desk, every single-name
mean-reversion signal is fitted by an overnight walk-forward OU
pipeline owned end-to-end by the Strats group. A typical morning
workflow on name \texttt{XYZ}:
\begin{enumerate}
  \item \textbf{Stationarity gate.} Run an ADF test on the
  fitting window. If the test fails to reject the random-walk
  null, the name does not trade today --- the $z$-score
  signal is meaningless on a non-stationary series.
  \item \textbf{Parameter estimation.} Fit
  $(\hat\theta, \hat\mu, \hat\volat)$ on the fitting window
  via the AR(1) mapping of \cref{subsec:ch9_ols}.
  \item \textbf{Half-life sanity check.} Compute
  $\hat h_{1/2} = (\ln 2)/\hat\theta$ and confirm it falls in
  an asset-class plausible range (intraday: minutes to hours;
  multi-day: days to weeks). Out-of-range flags the Strat to
  inspect.
  \item \textbf{$z$-score signal.} Stream live
  $z_t = (X_t - \hat\mu)/\hat\sigma_\infty$ to the trader's
  monitor and the automated execution layer.
  \item \textbf{Kelly sizing.} Convert $z_t$ into position size
  via Kelly (\cref{ch:kelly_risk}), capped at the desk's
  per-name VaR contribution.
\end{enumerate}
A standard GS/quant interview question:
\emph{``Your OU fit reports a half-life of $0.5$ days on this
name, but the average holding period is $3$ days. What do you
investigate?''} Three buckets in order: (i) a stale or
mis-applied $\hat\mu$ that has drifted from the live mean, so
$z_t$ takes longer to revert because the supposed mean is not
where the series is heading; (ii) a regime change that has
lengthened true $\theta$ since the fit, often visible by
re-fitting on the most recent sub-window; (iii) overfitting,
where $\hat\theta$ from a short window is biased toward larger
values than the truth. All three are diagnosable; walking
through them in order signals familiarity.
\Cref{ch:strats_role} returns to this kind of diagnostic
conversation.
\end{gsbox}

\section{Limitations and when OU breaks}
\label{sec:ch9_limits}

The OU model is the simplest non-trivial mean-reverter and the
default tool of practitioners. Like every default, it has
failure modes worth knowing before pointing it at production
money.

\paragraph{The mean is not constant.} OU assumes a \emph{single}
long-run mean $\mu$. Real series typically have a slow-moving
level on top of the fast-reverting noise --- for an equity, the
company's growing earnings; for an FX rate, the inflation
differential; for a Prosperity ETF premium, whatever the
simulator calibrated. Fix: de-trend before fitting (subtract a
moving average or a regression against a slow factor) or refit
$\hat\mu$ on a rolling window short enough that the drift inside
is small. Walk-forward handles this implicitly when the refit
cadence is fast enough.

\paragraph{The half-life is not constant.} OU assumes a single
$\theta$. Real $\theta$ varies with regime: faster in calm
markets, slower in stressed ones. Walk-forward picks this up to
first order, but a calm-regime strategy quietly stops working
when the half-life lengthens beyond the holding window.

\paragraph{The series is not stationary in level.} Many
``mean-reverting'' series are stationary only after
transformation: log, first difference, or residual against a
slow factor. Classic case: two cointegrated stocks --- each
price is a random walk, but their difference (with the right
hedge ratio) is OU. Trading the OU residual is
\cref{ch:pairs_cointegration}; the technique here applies
directly once cointegration is established. A related worry is
\emph{over-differencing}: aggressive differencing destroys the
long-run information carrying the signal. \citeAFML{} (Ch.~5)
treats this via \emph{fractional differentiation}. Rule of
thumb: use the \emph{lowest} order of transformation that makes
the series stationary.

\paragraph{Jumps and fat tails.} OU has Gaussian increments.
Real series jump (news, earnings, central-bank decisions, or
Prosperity's Olivia informed-trading windows) and the Gaussian
envelope underestimates the tails. A live OU strategy should
always layer a stop-out ($|z| > z_{\text{stop}}$, per
\cref{sec:ch9_zscore}) to bail out before a jump turns a signal
into a runaway position.

\paragraph{Cointegration without OU.}
\Cref{ch:pairs_cointegration} covers the next layer: many
tradeable opportunities live in \emph{linear combinations} of
series, where components are random walks but the combination
reverts. OU applies as soon as the combination is identified;
finding it, picking the hedge ratio, and testing it are the next
chapter.

\section{Key facts to memorise}
\label{sec:ch9_keyfacts}

\begin{itemize}
  \item The Ornstein--Uhlenbeck SDE is
  $dX_t = \theta(\mu - X_t)\,dt + \volat\,dW_t$, with
  $\theta > 0$, $\mu \in \R$, $\volat > 0$. It has a unique
  stationary distribution
  $\Normal(\mu,\, \volat^2/(2\theta))$.
  \item The closed-form solution
  $X_t = X_0 e^{-\theta t} + \mu(1 - e^{-\theta t})
       + \volat \int_0^t e^{-\theta(t-s)}\, dW_s$
  is obtained by multiplying the SDE through by the integrating
  factor $e^{\theta t}$ and integrating.
  \item \textbf{Half-life.} $h_{1/2} = (\ln 2)/\theta$. The
  single most important number for trading. If
  $h_{1/2} \ll H$ (your trading horizon), the strategy is
  viable; otherwise it is not.
  \item \textbf{Stationary standard deviation.}
  $\sigma_\infty = \volat/\sqrt{2\theta}$. The natural unit of
  the $z$-score.
  \item \textbf{Autocorrelation.}
  $\Corr(X_t, X_{t+h}) = e^{-\theta h}$. A single exponential
  with the same rate as the drift.
  \item \textbf{Discrete-time mapping.} OU sampled at spacing
  $\Delta t$ is exactly an AR(1) with $\phi = e^{-\theta \Delta t}$,
  intercept $\mu(1 - \phi)$, and innovation variance
  $\volat^2(1 - \phi^2)/(2\theta)$. Fit by OLS, recover
  $\hat\theta = -\ln(\hat\phi)/\Delta t$.
  \item \textbf{Stationarity tests.} ADF (regress
  $\Delta X_t$ on $X_{t-1}$ plus lags), variance ratio (compare
  $\Var(X_t - X_{t-k})$ against $k\,\Var(X_t - X_{t-1})$), and
  Hurst exponent ($H < 0.5$ mean reverting). Run all three;
  ADF alone has low power on short samples or slow reverters.
  \item \textbf{$z$-score rule.} Enter when $|z| > z_{\text{entry}}$,
  exit when $|z| < z_{\text{exit}}$, stop out when
  $|z| > z_{\text{stop}}$. Asymmetry between entry and exit
  thresholds is essential. Expected trade duration is
  $h \approx \theta^{-1}\ln(z_{\text{entry}}/z_{\text{exit}})$.
  \item \textbf{Walk-forward fitting is non-negotiable.} Fitting
  $\hat\mu$ once on the entire dataset and then using it for
  every historical $z$-score is the canonical look-ahead bias.
  Always refit on a rolling window using only data strictly
  before the trade time.
  \item OU is a \emph{variance-control} mean-reverter with
  Gaussian increments. It cannot model fat tails, jumps, or
  regime changes by itself --- always pair with a hard stop-out
  for protection.
  \item On Prosperity, KELP is mildly mean-reverting via a
  rounding artifact and is best traded with a fixed-offset
  market-maker (Frankfurt Hedgehogs); the OU $z$-score
  machinery shines on the PICNIC\_BASKET premium of
  \cref{ch:basket_etf_arb}, where half-lives are short and
  excursions large.
  \item On a sell-side stat-arb desk every single-name
  reversion signal is fitted by walk-forward OU; the Strat owns
  the pipeline ADF gate $\to$ estimation $\to$ half-life check
  $\to$ live $z$-score $\to$ Kelly sizing
  (\cref{ch:kelly_risk}). Standard interview probe:
  ``half-life says $X$, average holding is $Y$, what do you
  investigate?'' --- three buckets: stale $\hat\mu$, regime
  change, walk-forward overfitting.
\end{itemize}
